\scanpage{56}

il faut simplement voir que ceci ne dépend pas du choix de $\omega^{N-1}$, i.e. que \deleted{$d\omega^{N-1}=0$}
\[
d\omega^{N-1}=0\quad\Longrightarrow\quad
k_{N-1}(\partial\omega^{N-1})=0.
\]
Distinguons deux cas :

a) $N\geq2$ i.e. $N-1\geq1$, alors $H^{N-1}(A_N^*)=0$ par (1), donc
\[
d\omega^{N-1}=0\quad\Rightarrow\quad
\omega^{N-1}=d\alpha^{N-2},
\]
et on aura
\[
k_{N-1}(\partial\omega^{N-1})
=k_{N-1}(\partial d\alpha^{N-2})\unsure{\Leftarrow}.
\]
Or \deleted{$\unsure{d\omega^N}$} $\partial d\alpha^{N-2}=d\partial\alpha^{N-2}$ donc on \unsure{trouve}
\[
k_{N-1}(\partial\omega^{N-1})
=k_{N-1}(d\partial\alpha^{N-2})
\overset{\text{Stokes}}{=}
k_{N-2}(\underbrace{\partial\partial\alpha^{N-2}}_{0})=0
\]
cqfd

b) $N=1$ i.e. $N-1=0$. Alors en vertu de (1) on
$d\omega^0=0$ implique $\omega^0=\varepsilon_1(\alpha)$
$(\alpha\in A_0)$ donc
\[
k_{N-1}(\partial\omega^0)
=k_{N-1}(\underbrace{\partial\varepsilon_1(\alpha)}_{\text{''}0})=0
\qquad\deleted{\unsure{0=\partial\varepsilon_1(\alpha)}}\quad\text{o.k.}
\]

\uline{Remarques.} 1) Dans \inserted{l'hyp.} (1), en plus de la condition des degrés, on a utilisé
\[
\left\{
\begin{aligned}
H^N(A_N^*)&=0 &&\text{si }N\geq1\\
H^{N-1}(A_N^*)&=
\begin{cases}
0&\text{si }N\geq2\\
\varepsilon_1(A_0)&\text{si }N=1
\end{cases}
\end{aligned}
\right.
\]

2) On peut donner une version ``abstraite'' en termes d'un bicomplexe augmenté avec des conditions de degré (\unsure{qu'il faut} (et $d\partial=\partial d$ \unsure{!})) \deleted{\unsure{degré}} dans une cat. abélienne. [\unsure{ceci peut être} $(\mathrm{Ab})^{\unsure{\circ}}$ !]

\diagram{\begin{tikzcd}[column sep=small,row sep=small]
&&&0&A_3^3&\arrow[l,dash]\\
&&0&A_2^2\arrow[l]\arrow[u]&A_3^2\arrow[l,"\partial"']\arrow[u,"d"]&\arrow[l]\\
&0&A_1^1\arrow[l]\arrow[u]&A_2^1\arrow[l,"\partial"']\arrow[u,"d"]&A_3^1\arrow[l,"\partial"']\arrow[u,"d"]&\arrow[l,dash]\\
0&A_0\arrow[l]\arrow[u]&A_1^0\arrow[l,"\partial"']\arrow[u,"d"]&A_2^0\arrow[l,"\partial"']\arrow[u,"d"]&A_3^0\arrow[l,"\partial"']\arrow[u,"d"]&\arrow[l,dashed]\\
0\arrow[r]&A_0\arrow[u,"\varepsilon_0","\mathrm{id}"']&A_0\arrow[l,"0"']\arrow[u,"\varepsilon_1"]&A_0\arrow[l,"\mathrm{id}"']\arrow[u,"\varepsilon_2"]&A_0\arrow[l,"0"']\arrow[u,"\varepsilon_3"]&\arrow[l,"\mathrm{id}"']
\end{tikzcd}}

Conditions \deleted{$\unsure{A_N^i=0}\quad H^N(A_N^*)=0\text{ si }N\geq1$}
\[
\left\{
\begin{aligned}
A_N^i&=0&&\text{si }i\notin[0,N]\\
H^N(A_N^*)&=0&&\text{si }N\geq1\\
H^{N-1}(A_N^*)&\ \deleted{\unsure{=}}\ 
\begin{cases}
=0&\text{si }N\geq2\\
\subset\operatorname{Ker}\partial&\text{si }N=1
\end{cases}
\end{aligned}
\right.
\]

3) À cause de l'unicité, \unsure{ou} grâce à la version ``cat. abélienne'', on voit que les hom ``intégration''
$A_n^n\longrightarrow A_0^0=A_0$ commutent à tous les endom des $A_*^*$, donc respectent \unsure{les éventuelles} structures de $k$-modules.

\margin{55}
